La fuente formalizada es el manuscrito propio fijado en el commit
\href{https://github.com/sbvasquezg-ux/ai-project/commit/3fb4aab2a2064fb39fa79d48bcfcbc49eb7920fd}{\texttt{3fb4aab}}.
Los enunciados, las pruebas matemáticas y sus condiciones se conservan en la entrega del 15 de noviembre; solo se actualizan presentación, secuencia del período, declaración de colaboración y alcance manuscrito. Por ello la correspondencia de los dos resultados con la fuente fijada no cambia. Este apéndice informa su formalización. La carpeta se desarrolló en AppliedModelingLib, commit \texttt{194edcc}, y se exporta completa a \texttt{lean/}, respetando su \texttt{.gitignore}.
\begin{center}\small
\begin{tabular}{>{\raggedright\arraybackslash}p{.16\linewidth}>{\raggedright\arraybackslash}p{.35\linewidth}>{\raggedright\arraybackslash}p{.39\linewidth}}\toprule
Resultado & Declaración / archivo & Estado y alcance\\\midrule
Prop. 1 & \texttt{proposition\_one} en \texttt{ProofInterface.lean}; \texttt{PropositionOneSpec} & Probada: FOC, cociente $hr$, identidad con $r=0$ y pendiente a partir de beneficio cero en un entorno.\\[5pt]
Prop. 2 & \texttt{proposition\_two} en \texttt{ProofInterface.lean}; \texttt{PropositionTwoSpec} & Probada: existencia y unicidad de asignaciones y salarios en ambas economías; renta cero, elegibilidad y ganancia neta.\\\bottomrule
\end{tabular}
\end{center}
La definición de equilibrio representa las nueve actividades, los recursos humanos con tiempo de ayuda, el cómputo, la ausencia de beneficio y la complementariedad. La renta $r=0$ se deduce de $\mu>h$; no se introduce como una conclusión asumida. Las condiciones del régimen son las primitivas y desigualdades de la Proposición 2. No se añade una hipótesis de optimalidad, existencia o unicidad.

La prueba formal de unicidad usa la identidad de brecha salarial, las desigualdades de ambas economías y las holguras estrictas para descartar actividades alternativas. Los recursos fijan las masas restantes y la complementariedad fija salarios y renta. El resultado es global dentro del régimen de tres tipos; no es un teorema del continuo.

Los registros reproducibles se conservan en \texttt{lean/}. El punto de entrada del curso es:
\begin{quote}\small\ttfamily
python3 scripts/paper\_contribution.py check\\
\hspace*{1em}Vasquez26HeterogeneousVerification --fast
\end{quote}
Las declaraciones están cerradas sin huecos ni axiomas añadidos. La inspección de ambas pruebas usa únicamente los axiomas estándar \texttt{propext}, \texttt{Classical.choice} y \texttt{Quot.sound}. La autora declaró su aceptación de las contribuciones el 15 de noviembre. Esa declaración no sustituye evidencia manuscrita adicional: solo se entrega la hoja sobre la cota de $\kappa$, como se explicita en el último apéndice.

La corrida se ejecutó con GPT-5.6 Sol/xhigh sobre repositorios públicos. Tras sus ajustes y compilación, se agotó la cuota durante una auditoría adicional. El asistente coordinador ejecutó los chequeos finales, ambos con salida 0, y un revisor independiente comparó el código final con las dos proposiciones. El informe y los límites del registro están en \texttt{lean/docs/COURSE\_VALIDATION\_REPORT.md}. No se reivindica un recibo avanzado de publicación de AppliedModelingLib ni una auditoría técnica independiente realizada por la autora. Su aceptación declarada y la comprobación manuscrita aportada se registran por separado en el apéndice de colaboración.
